\subsection{Higher smooth-cap obstruction for the OS49 source}
\begin{statement}{P3K.81}{Proved}
For the fixed OS49 source and its component Cartier line, every
actual sphere member obtained by the stated finite root caps and smooth
translation caps has only order-one effective caps at smooth source values.
This holds with the full norm vectors, all finite characters and other
specialization groups retained, and for every allowed lattice-basis cusp
fan. At any effective smooth cap of order \(r>1\), the actual degree-two
specialization quotient is \((\mathbb Z/r)^3\). The unique global invariant
\(\omega=uw+3\psi\delta\) leaves a nonzero group in
\(H^1(\mathbb P^1,R^2\pi_*\mathbb Z)\), which survives in \(H^3(X;\mathbb Z)\).

For the explicit retained finite cap \((m,n,d)=(3,-1,0)\) and one smooth
cap of order two with full norm vector \(\gamma=(1,0,0,0)\), the actual
constructed compact manifold has
\[
 \pi_1(X)=1,\qquad H_2(X;\mathbb Z)=H_3(X;\mathbb Z)=(\mathbb Z/2)^2,
\]
\[
 H_0(X;\mathbb Z)=H_6(X;\mathbb Z)=\mathbb Z,
 \qquad H_1(X;\mathbb Z)=H_4(X;\mathbb Z)=H_5(X;\mathbb Z)=0.
\]
The proof supplies the full specialization, integral Mayer--Vietoris,
translation-prism and fundamental-group maps.

At the type-II value, every realized free good root has effective order
divisible by six. Its actual root source and corrected Cartier line give
\(\psi(\ell)=n\) for the full norm vector. Together with the retained
\(IV^*\) cap of order \(3k\), it has the displayed nonzero
\(\mathbb Z/3\) fundamental-group quotient and cannot be a sphere.
The reduced-II orders \(6l+1\) and their return classes are not assigned
by this statement.
\end{statement}
\begin{proof}
The following full calculation has passed the recorded independent audit. Retain the actual OS49 source and
component line, the finite IV-star cap \((m,n,d)=(3,-1,0)\), and a fixed smooth
source value c.~Add the ordinary cap of order two whose full norm vector
is \(\gamma=(1,0,0,0)\) in the actual \((\gamma,u,w,\delta)\) homology marking. No
return class is assigned. The completed manifold is simply connected,
but its integral groups are
\[
 H_2(X)=H_3(X)=(\mathbb Z/2)^2,
 \quad H_1(X)=H_4(X)=H_5(X)=0,
 \quad H_0(X)=H_6(X)=\mathbb Z.                     \tag{OT.1}
\]
All maps below are the actual local specialization and integral
Cech/Mayer--Vietoris maps, with the original four fibre directions retained.

\par

\medskip\noindent\textbf{OT.1. The actual local cap and its pullback in every degree}

\par\nobreak

On a sufficiently small disc about c, the source and its degree-zero line
are smooth. Retain the actual source translation and its complete line
unit in the suspension B. On this disc choose their holomorphic Lie
logarithms. Half of this actual logarithm gives a commuting lift with
square B, including its line multiplier. On the root \(q=z^2\) the compact
torus action is
\[
                  (x,z)\longmapsto
                 (x+\Pi(q)\gamma/2,-z).                 \tag{OT.2}
\]
Its square adds the actual integral B period. Its nontrivial element is
free because the integral covector \(\psi\) evaluates \(\gamma\) to one. The
ordinary-to-root overlap is the actual plus-log translation with \(\gamma\),
so its filled meridian is \(2\alpha_S-\gamma\). All source and line units used
to form the half logarithm remain in this overlap.

The cap retracts onto the compact torus with period lattice
\(\Lambda'=\Lambda+\mathbb Z(\gamma/2)\). In the actual dual basis \((\psi,u,w,\delta)\), its
first-cohomology pullback is exactly
\[
                R_1=\operatorname{diag}(2,1,1,1),
                 \qquad R_q=\bigwedge^q R_1.              \tag{OT.3}
\]
This follows by evaluating an integral covector on \(\gamma\)/2; it is not a
rational rank computation. In lexicographic exterior bases \(R_q\) has
entry two precisely on basis forms containing \(\psi\), and entry one on
all the others. Thus the specialization quotient at this point is
\[
           Q_{S,q}=(\mathbb Z/2)^{\binom3{q-1}}
                 \quad(1\le q\le4).                    \tag{OT.4}
\]
The quotient map is the literal vector of coefficients on the basis
forms containing \(\psi\), reduced modulo two.

At the retained order-three cap, the norm is -\(\ell_0\) with
\(\ell_0=(1,-1,0,0)\). In actual root coordinates the map \(x\mapsto-x\) conjugates
the affine translation \(+\ell_0/3\) to \(-\ell_0/3\), keeps its linear action A,
and carries every full fibre period to its inverse. It extends on this
smooth root torus family. Hence the actual specialization lattices are
those computed in F26, with the sign \((-1)^q\) retained on degree q. Their
quotients are \(\mathbb Z/3\) in degrees one and four, and zero in degrees zero,
two and three. More explicitly the invariant degree-one basis is
\(\psi,\ h=-w+3\delta\), and the cap image is generated by \(3\psi,h\). In degree two
the full invariant lattice generated by
\[
                u_2=-\psi w+3\psi\delta,
                 \qquad v_2=\psi w+uw
\]
is its image. Degree three has the full invariant image, and degree four
has image \(3\nu\). These are integral cap maps from the actual finite
quotient, including its norm sign; no cohomology torsion is omitted.
The retained II and selected component cusp have the full integral
invariant specialization images of F26.

\par

\medskip\noindent\textbf{OT.2. The complete sheaf Mayer--Vietoris presentation}

\par\nobreak

Write \(L_q\) for the degree-q fibre local system and \(J_q=j_*L_q\) for its
invariant extension across the three original critical values. The
actual cohomological monodromies are
\[
 T_A=\begin{pmatrix}1&-3&-3&-1\\0&-2&-3&-1\\0&1&1&0\\0&0&0&1\end{pmatrix},
 \quad T_B=\begin{pmatrix}1&0&0&0\\0&1&1&0\\0&-1&0&0\\0&0&0&1\end{pmatrix},
 \quad T_C=\begin{pmatrix}1&0&3&1\\0&1&2&1\\0&0&1&0\\0&0&0&1\end{pmatrix}.
                                                               \tag{OT.5}
\]
The smooth new puncture has identity linear action. For \(E_i=\bigwedge^q T_i\),
the integral parabolic complex is
\[
 L_q\xrightarrow{d^0}\bigoplus_i(E_i-I)L_q
       \xrightarrow{d^1}L_q,
 \quad d^0v=((E_i-I)v)_i,
 \quad d^1(y_A,y_B,y_C)=y_A+E_Ay_B+E_AE_By_C.             \tag{OT.6}
\]
The exact Hermite bases, both maps and their full Smith forms are retained
in the certificate. Their middle cohomology is zero integrally, and the
last cokernel is free of rank one. The generators and projections are
\[
\begin{array}{c|c|c}
q&H^0(J_q)&H^2(J_q)/H_B\\\hline
0&1&[1]\\
1&\psi&[\delta]\\
2&\omega=uw+3\psi\delta&\eta,\quad[\psi\delta]=\eta,\ [uw]=3\eta\\
3&\phi=\psi uw&[uw\delta]\\
4&\nu=\psi uw\delta&[\nu].
\end{array}                                                   \tag{OT.7}
\]
Every other exterior-basis coefficient maps to zero in the displayed
coinvariant projection. All these statements use the full rank-four
matrices; the component line has not been replaced by a whole-fibre line.

Let \(S_q=R^q\pi_*\mathbb Z\) for the actual two-cap family. Every local restriction
above is injective, including the cap torsion-free cohomology, so the
literal restriction maps give
\[
 0\longrightarrow S_q\longrightarrow J_q
   \longrightarrow Q_{A,q}\oplus Q_{S,q}\longrightarrow0.       \tag{OT.8}
\]
The last map is supported at the two cap points. On a cover by the
ordinary complement and their discs, (OT.8) is precisely the quotient of
the actual integral Cech/Mayer--Vietoris restriction complex by the local
invariant blocks; its finite indices are (OT.3) and the order-three maps
above. It retains all base overlaps. In particular it is not a substitution
of invariant rational cohomology for the finite cap.

The global-invariant map to the skyscraper group has the following full
integer presentations, with the last column being that map:
\[
\begin{array}{c|c|c}
q&[\text{local relations}\mid\text{global image}]&H^1(S_q)\\\hline
1&\begin{pmatrix}3&0&1\\0&2&1\end{pmatrix}&0\\[3pt]
2&\begin{pmatrix}2&0&0&0\\0&2&0&0\\0&0&2&3\end{pmatrix}&(\mathbb Z/2)^2\\[3pt]
3&\begin{pmatrix}2&0&0&1\\0&2&0&0\\0&0&2&0\end{pmatrix}&(\mathbb Z/2)^2\\[3pt]
4&\begin{pmatrix}3&0&1\\0&2&1\end{pmatrix}&0.
\end{array}                                                    \tag{OT.9}
\]
For \(q=2\) the three smooth quotient coordinates are \(\psi u,\psi w,\psi\delta\), and the global image is \((0,0,3)\). For \(q=3\) they are \(\psi uw,\psi u\delta,\psi w\delta\), and the image is \((1,0,0)\). Thus these are the
actual complete maps, rather than only their Smith factors.

The kernels give the primitive zeroth-column basis
\[
                    1,\quad6\psi,\quad2\omega,
                          \quad2\phi,\quad6\nu.          \tag{OT.10}
\]
The two-cap sheaf still has the same free second column as (OT.7), because
the supported quotient in (OT.8) has no positive cohomology. There are no
other columns. Formula (OT.9) displays the nonzero middle column that a
calculation using only the Seifert determinant would miss.

\par

\medskip\noindent\textbf{OT.3. Actual translation prisms and the four integral differentials}

\par\nobreak

The untwisted source origin is the one constructed in F26: a meromorphic
frame of M restricted to O has its divisor at the selected cusp; the
actual integral cusp period is subtracted, so the origin extends there
and at II. Add the new smooth disc with its unchanged zero origin.
For the auxiliary linear root quotients, inversion about these origins
commutes with their linear monodromies and extends on the centrally
symmetric same toric fan. Its action on fibre degree q is \((-1)^q\).
Naturality therefore makes the intrinsic rational \(d_2\) zero. This is the
same actual origin and full line/cusp comparison as F26, with one
additional ordinary zero chart.

The actual affine finite-cap replacement is compared on its boundary by
the retained translation prism. Three lifted copies of the first prism
close to the integral translation -\(\ell_0\); two copies of the new prism
close to \(\gamma\). Their full operators on a fibre exterior form are
therefore
\[
             d_2=-\frac13\iota_{\ell_0}
                         +\frac12\iota_\gamma.           \tag{OT.11}
\]
The line and source entries of \(\ell_0\), including its -u entry, remain in
this formula. The source/target of each actual \(d_2\) from column zero to
column two is a free integral rank-one lattice, so its rational value
on the primitive bases (OT.7),(OT.10) determines its integral value.
The four evaluations are all exactly one:
\[
 d_2(6\psi)=H_B,\quad
 d_2(2\omega)=H_B[\delta],\quad
 d_2(2\phi)=H_B\eta,\quad
 d_2(6\nu)=H_B[uw\delta].                              \tag{OT.12}
\]
For example the second raw contraction is \(\delta+(2/3)w\); its w term is
zero under the actual coinvariant projection. The third is
\((uw-2\psi w)/3\), whose projection is \(\eta\). The fourth is
\(uw\delta-2\psi w\delta\), with the displayed target. These checks retain
all coordinates before projection.

The middle-column groups in (OT.9) have no possible incoming or outgoing
differential: the base has cohomological dimension two. All free terms
except degrees zero and six cancel by (OT.12). Hence
\[
 H^3(X;\mathbb Z)=H^4(X;\mathbb Z)=(\mathbb Z/2)^2,
 \qquad H^1=H^2=H^5=0.                                  \tag{OT.13}
\]
There are no extension ambiguities because each of these middle degrees
has only the single displayed surviving filtration group.

\par

\medskip\noindent\textbf{OT.4. Full fundamental group and integral homology}

\par\nobreak

Use the actual homological matrices \(H_i=(T_i^{-1})^t\). The II and cusp
fillings kill their base meridians and the images of \(H_B-I,H_C-I\). This
rank-three primitive relation module has quotient \(\mathbb Z\gamma\), as in F26.
The remaining base loops are \(\alpha_A,\alpha_S\). The exact finite-cap and
base relations reduce to
\[
 3\alpha_A+\gamma=0,\qquad 2\alpha_S-\gamma=0,
                  \qquad \alpha_A+\alpha_S=0.             \tag{OT.14}
\]
Their matrix on \((\gamma,\alpha_A,\alpha_S)\) is
\[
                  \begin{pmatrix}1&3&0\\-1&0&2\\0&1&1\end{pmatrix},
                     \qquad\det=1.                      \tag{OT.15}
\]
These are the full van Kampen relations, not merely a presentation of an
abelianization: the original torus generators commute, the killed II/cusp
meridians impose the displayed monodromy relations, and the base relation
leaves a single base generator. The surviving invariant \(\gamma\) commutes
with it. Every boundary maps surjectively onto its local filled group.
The full six-generator, eleven-column integer matrix and the actual
preimages of the \(H_A-I\) columns are retained in the certificate. Thus
\(\pi_1(X)=1\). Universal coefficients and integral Poincare duality applied
to (OT.13) give exactly (OT.1).

The chosen local cap is smooth, since its nontrivial affine deck is free;
the original root and cusp charts are unchanged. Compactness follows by
the same actual compact fundamental-domain gluing as F26. The calculation
therefore concerns the specified constructed manifold, not only a formal
Seifert presentation.

\par

\medskip\noindent\textbf{OT.5. The general primary obstruction for added smooth caps}

\par\nobreak

Retain any actual free OS49 finite cap of order \(m_0=3k\), and add finitely
many effective smooth logarithmic caps of orders \(m_i\) with their complete
norm vectors \(\ell_i\). The invariant covector \(\psi\) measures \(n_i=\psi(\ell_i)\).
Quotient the actual fundamental-group presentation by the normal fibre
subgroup killed by \(\psi\). The remaining marked relations are
\[
 m_i\alpha_i-n_i\gamma=0,\qquad
                  \sum_i\alpha_i+b_0\gamma=0.             \tag{OT.16}
\]
Here \(b_0\) retains all ordinary order-one clutches and the actual base
framing. It is not set to zero.

A sphere member must have pairwise coprime effective orders. Indeed if a
prime p divides two orders, all other invertible-order meridians can be
eliminated modulo p, while the indicated two have only their one sum
relation; their own norm equations constrain only \(\gamma\). This quotient
has a nonzero first homology class over \(\mathbb F_p\). When only \(m_i\) is divisible
by p, its coefficient \(n_i\) must be a unit modulo p, for otherwise \(\gamma\)
remains unconstrained after the same elimination. These are exact
consequences of (OT.16), independent of the other entries of \(\ell_i\).

Fix a prime p dividing the order of an added smooth cap in a sphere datum,
and use \(R=\mathbb Z_{(p)}\). Then p does not divide \(m_0\), in particular p is not three,
and no other effective order is divisible by p.~All the other finite
cap pullbacks become isomorphisms onto their invariant R lattices by
the actual finite covering and transfer; their degrees are units in R.
The II and cusp maps remain the integral invariant maps. The unchanged
parabolic complex (OT.6) has \(H^1=0\) over R and invariant degree-two
generator \(\omega\).

Write the full chosen norm vector as \(\ell=(n,a,b,d)\), with n a unit of R.
An explicit basis of its annihilator in the dual lattice is
\[
 u-(a/n)\psi,\qquad w-(b/n)\psi,
                         \qquad\delta-(d/n)\psi.          \tag{OT.17}
\]
The local degree-two quotient is \((R/m_iR)^3\), by the same actual torus
covering as (OT.3), now with \(\Lambda+R(\ell/m_i)\). Its natural quotient map
is contraction with \(\ell\) modulo \(m_i\). On the retained global invariant,
\[
 \iota_\ell\omega=-b u+a w+3n\delta-3d\psi,
 \qquad\text{coordinates in (OT.17): }(-b,a,3n).           \tag{OT.18}
\]
The third coefficient is a unit in R. The full local/global MV
presentation is therefore
\[
            [\ m_i I_3\mid(-b,a,3n)^t\ ],
 \qquad H^1(R^2\pi_*R)=(R/m_iR)^2\ne0.                 \tag{OT.19}
\]
The source, line and \(\delta\) coefficients are retained in (OT.17)--(OT.18);
this is a calculation in the localized lattice, not a change of the
actual cap or a deletion of its \(\delta\) character.

This middle-column group survives to the integral localized cohomology
filtration for the same dimensional reason as in OT.3. Thus \(H^3(X;R)\)
has a nonzero filtration quotient, contradicting the sphere cohomology.
Consequently every added effective smooth cap of an OS49 sphere member
has order one. The actual periodic flop comparisons of F35 are the
identity near every cap and preserve integral homology and \(\pi_1\), so the
obstruction applies to every allowed lattice-basis fan without using an
unproved return formula.

The executed certificate is \path{code/part3/os49_order_two_cap.py}, with all
exterior matrices, parabolic maps, local quotient presentations, primitive
indices, prism contractions and the full fundamental relation matrix.
Its first run passes 63 identities. The preceding localized argument is
retained in full. The next section gives a stronger direct obstruction
that does not use the pairwise-coprime reduction or transfer at the other
caps.

\par

\medskip\noindent\textbf{OT.6. Every higher effective smooth cap is excluded directly}

\par\nobreak

Retain the fixed OS49 source, its actual component line and linear
period tuple (OT.5). Take any realized free finite caps at the other
sites and any finite collection of actual smooth logarithmic caps. The
word effective records the order of the translation on the central
compact torus. On a smooth root disc the full actual affine map has the
form
\[
 (x,z)\longmapsto(x+\Pi(z^r)\ell/r,\zeta_r z),
 \qquad \ell\in\Lambda,\qquad
       \gcd(r,\ell_1,\ell_2,\ell_3,\ell_4)=1.          \tag{OT.20}
\]
Here is the full map to these coordinates. In the original smooth root
atlas write the actual lifted translation as a(z), with its source and
line entries, and choose its holomorphic universal-cover lift on the
disc. The literal norm gives
\(\sum_{j=0}^{r-1}a(\zeta_r^jz)=\Pi(z^r)\ell\). Put
\[
 f(z)=a(z)-\Pi(z^r)\ell/r,\qquad
 g(z)=\frac1r\sum_{j=0}^{r-1}j f(\zeta_r^j z).
                                                               \tag{OT.20a}
\]
Then \(g(\zeta_rz)-g(z)=f(z)\), so \((x,z)\mapsto(x-g(z),z)\)
intertwines the original deck with (OT.20). This map commutes with the
full period lattice and the physical circle, extends on the root disc,
and is carried through the actual ordinary-to-cap overlap. Its constant
finite character is part of \(\ell\) and is not removed by this averaging.
The gcd condition says precisely that
the translation has order r, so every nonidentity deck element acts
freely at the central fibre. A chosen root with a smaller translation
order has a subgroup acting only on the root coordinate. Its quotient
is the literal map \(z\mapsto z^g\), where
\(g=\gcd(r,\ell_1,\ell_2,\ell_3,\ell_4)\); the residual displayed map has order
\(r/g\) and vector \(\ell/g\). Thus (OT.20) specifies the effective cap without
changing an effective nontrivial translation into an order-one cap.

The actual cap retracts to
\(T'=\Lambda_{\mathbb R}/(\Lambda+\mathbb Z\ell/r)\), and the specialization from an ordinary
fibre is the identity-linear torus covering
\[
 c:\Lambda_{\mathbb R}/\Lambda\longrightarrow
       \Lambda_{\mathbb R}/(\Lambda+\mathbb Z\ell/r).
                                                               \tag{OT.21}
\]
In the real flat marking \(x=\Pi(q)y\), the radial retraction
\((y,z)\mapsto(y,tz)\), \(0\le t\le1\), commutes with the deck
\((y,z)\mapsto(y+\ell/r,\zeta_rz)\). It proves both the asserted cap retraction
and the actual specialization covering (OT.21). Its degree is r and
its deck group is generated by \(\ell/r\). In particular
the first-cohomology pullback is the actual sublattice
\[
 c^*H^1(T';\mathbb Z)
 =\{v\in\Lambda^*:v(\ell)\equiv0\pmod r\}
                 \subset\Lambda^*.                         \tag{OT.22}
\]
The evaluation map onto \(\mathbb Z/r\) is surjective by the gcd condition.
Consequently (OT.22) has Smith factors \((1,1,1,r)\). Exterior powers give
the literal degree-two specialization quotient
\[
 Q_{c,2}=\bigwedge^2\Lambda^*/
                 \bigwedge^2c^*H^1(T';\mathbb Z)
                      \simeq(\mathbb Z/r)^3.              \tag{OT.23}
\]
For an entirely integral coordinate construction, let H be the column
Hermite basis of \([rI_4\mid\ell]\). Then \(\Lambda'=r^{-1}H\mathbb Z^4\), the actual
dual pullback matrix is \(rH^{-t}\), and its exterior square is the map
in (OT.23). This retains every entry of \(\ell\) and furnishes the exact
lattice morphism underlying the Smith description.

Let \(R=R^2\pi_*\mathbb Z\) and \(J=j_*L_2\) on the fixed base. On the ordinary locus
the specialization map R to J is the identity. Let K be its kernel,
S its image, and Q its cokernel. Their actual stalk maps give
\[
 0\longrightarrow K\longrightarrow R\longrightarrow S
       \longrightarrow0,\qquad
 0\longrightarrow S\longrightarrow J\longrightarrow Q
       \longrightarrow0.                                \tag{OT.24}
\]
K and Q are supported on the finite set of special values. This retains
all kernel and torsion groups of every other cap or cusp completion;
they are not assumed to vanish. Such a finite-support sheaf has no
positive sheaf cohomology. Hence
\[
 H^1(R)\simeq H^1(S),\qquad
 \operatorname{coker}\{H^0(J)\longrightarrow H^0(Q)\}
                         \hookrightarrow H^1(S).         \tag{OT.25}
\]
The fixed tuple gives \(H^0(J)=\mathbb Z\omega\), as checked integrally in OT.2.
Projection of \(H^0(Q)\) to the actual smooth-cap summand (OT.23) induces
a surjection from the cokernel in (OT.25) onto
\[
                  Q_{c,2}/\langle\operatorname{sp}_c(\omega)\rangle.
                                                               \tag{OT.26}
\]
For every prime p dividing r, the quotient in (OT.26), tensored with
\(\mathbb F_p\), has dimension at least two: (OT.23) has three generators and the
displayed image supplies only one relation. It is therefore nonzero.
Equations (OT.25)--(OT.26) prove \(H^1(R^2\pi_*\mathbb Z)\) nonzero without any
condition on specialization at the other caps.

The Leray term \(E_2^{1,2}\) has no possible incoming or outgoing
differential. Incoming indices have negative base degree; outgoing
indices exceed the cohomological dimension two of the real base.
Thus it survives as a nonzero associated-graded part of \(H^3(X;\mathbb Z)\).
This excludes sphere cohomology for every effective smooth cap of
order r greater than one.

This also proves the asserted all-fan scope directly. On the dense
ordinary locus the actual toric changes are the identity with the same
full period tuple, and (OT.21) is unchanged near every smooth cap.
Changing the selected-cusp fan changes only a special stalk already
retained in K and Q. For the full homology calculation of OT.1, the
actual common blowups and inverse blowdowns of F35 give the integral
homology and fundamental-group transport, relative to all cap charts;
both exceptional rulings pair primitively with the exceptional divisor.
No return formula or identification of the two rulings is used here.

\par

\medskip\noindent\textbf{OT.7. The actual good-II root and its surviving norm}

\par\nobreak

At the source value \(t=1\) put \(q=t-1\). The exact OS49 equation is
\[
 y^2=x^3-3(1+q)^3q x+(1+q)^4q(1+2q),
                       \qquad G=((1+q)^2,(1+q)^2).       \tag{OT.27}
\]
For a good root \(m_0=6l\), l at least one, put \(q=z^{6l}\),
\(x=z^{2l}X\), \(y=z^{3l}Y\). Its minimal equation and marked section are
\[
 \begin{split}
 Y^2&=X^3-3(1+z^{6l})^3z^{2l}X
                   +(1+z^{6l})^4(1+2z^{6l}),\\
 X(G')&=z^{-2l}(1+z^{6l})^2,\qquad
 Y(G')=z^{-3l}(1+z^{6l})^2,\qquad
                  -X(G')/Y(G')=-z^l .
 \end{split}                                                \tag{OT.28}
\]
Thus \(G'\text{ and }Q'=3G'\) both tend to O on the central curve
\(Y^2=X^3+1\). The actual deck is
\[
 D:(z,X,Y)\longmapsto
       (\zeta_{m_0}z,\zeta_{m_0}^{-2l}X,
                              \zeta_{m_0}^{-3l}Y).        \tag{OT.29}
\]
Its central order is six. The small formal-group inverse of
\(m_0\) gives the D-invariant section R with \(m_0R=Q'\) and \(R(0)=O\).

The local six-point blowup/quotient/contraction diagram of F43 applies
to the two marked sections G' and O individually. Both meet the same
order-six exceptional curve at finite coordinates, whereas all three
contractions occur at its infinity point. Their exceptional coefficients
cancel in their difference. The literal Cartier identity is therefore
\[
 \beta^*\mathcal O(G'-O')
               =\pi^*\alpha^*\mathcal O(G-O).             \tag{OT.30}
\]
This is the local diagram of F43 with \(d=1\); its proof uses the actual two
section divisors and requires no change to the component-supported
line at the distant selected \(I_2\) fibre. Translation Q' tends to O and
fixes all six centres, so (OT.30) carries the actual source/line
suspension and its deck action through that diagram.

Lift \(t_{-R}\) on this actual root line to F. In the retained finite
character gauge, \(DFD^{-1}=\chi F\), \(\chi^{m_0}=1\), and
\(B_0=\beta(q)F^{m_0}\). The invariant norm-root correction from F43 gives
the commuting full lift \(A_0\) with \(A_0^{m_0}=B_0^{n_0}\); it retains the
factor \(\chi^{n_0m_0(m_0-1)/2}\) and the prescribed \(\delta\) character.
The coefficient covector \(\psi\) kills the elliptic and physical-\(\delta\)
periods and takes value one on the full suspension period B. Lifting
the literal root relation \(A^{m_0}=B^{n_0}\) to the actual lattice thus
gives
\[
 \psi(\ell_0)=n_0,\qquad
          m_0\alpha_0-\ell_0=0,
          \quad m_0\alpha_0-n_0\gamma=0
                       \text{ after the psi quotient}. \tag{OT.31}
\]
All other source and \(\delta\) entries of \(\ell_0\) remain in the first
relation. This establishes the actual norm before using the base
determinant.

Every such effective free good-II cap has order divisible by six, and
the retained IV-star cap has order 3k. There is an explicit surjection
of the full fundamental-group presentation onto \(\mathbb Z/3\): send the
IV-star base meridian to 1, the good-II meridian to -1, every other
base meridian to zero, and every fibre generator to zero. Both
multiple-cap relations vanish because their orders are divisible by
three; the base sum and all affine conjugation relations also vanish.
The local free quotient maps retain exactly these root meridians, so
this is a homomorphism on the filled manifold. Hence a good-II root
replacement cannot be a sphere member. The reduced-II orders
\(6l+1\) are not assigned a topology or return by this calculation.

\end{proof}
